%% ****** Start of file apstemplate.tex ****** %
%%
%%
%%   This file is part of the APS files in the REVTeX 4.2 distribution.
%%   Version 4.2a of REVTeX, January, 2015
%%
%%
%%   Copyright (c) 2015 The American Physical Society.
%%
%%   See the REVTeX 4 README file for restrictions and more information.
%%
%
% This is a template for producing manuscripts for use with REVTEX 4.2
% Copy this file to another name and then work on that file.
% That way, you always have this original template file to use.
%
% Group addresses by affiliation; use superscriptaddress for long
% author lists, or if there are many overlapping affiliations.
% For Phys. Rev. appearance, change preprint to twocolumn.
% Choose pra, prb, prc, prd, pre, prl, prstab, prstper, or rmp for journal
%  Add 'draft' option to mark overfull boxes with black boxes
%  Add 'showkeys' option to make keywords appear
%\documentclass[aps,prl,preprint,groupedaddress]{revtex4-2}
\documentclass[aps,prl,twocolumn,groupedaddress]{revtex4-2}
%\documentclass[aps,prl,preprint,superscriptaddress]{revtex4-2}
%\documentclass[aps,prl,reprint,groupedaddress]{revtex4-2}

% You should use BibTeX and apsrev.bst for references
% Choosing a journal automatically selects the correct APS
% BibTeX style file (bst file), so only uncomment the line
% below if necessary.
%\bibliographystyle{apsrev4-2}

\usepackage{graphicx}
\usepackage{epstopdf}
\usepackage{caption}
%\usepackage{amsmath}% http://ctan.org/pkg/amsmath
%\usepackage{amsthm}
%\usepackage{amsfonts}
%\usepackage{subfigure}
%\usepackage{hhline}
%\usepackage[miktex]{gnuplottex}
%\usepackage{xcolor}
\usepackage{amssymb}
\usepackage{amsmath}
\usepackage{color}
\usepackage{hyperref}
%\usepackage[percent]{overpic}
\usepackage{tikz}
\usepackage{mathrsfs}
\usepackage{wasysym}
\usepackage{tikz-cd}
%\usepackage{linearb}
%\usepackage{wsuipa}
\usepackage{calc}
\usepackage{tipa}

%\usepackage{stix} %\fisheye
\usepackage{stackengine,scalerel}

% so sections, subsections, etc. become numerated.
\setcounter{secnumdepth}{3}

\newcommand{\FF}{\mathbb{F}}
\newcommand{\NN}{\mathbb{N}}
\newcommand{\ZZ}{\mathbb{Z}}
\newcommand{\QQ}{\mathbb{Q}}
\newcommand{\RR}{\mathbb{R}}
\newcommand{\CC}{\mathbb{C}}
\newcommand{\II}{\mathbb{I}}
\newcommand{\GG}{\mathbb{G}}

\DeclareMathOperator*{\argmax}{arg\,max}
\DeclareMathOperator*{\argmin}{arg\,min}
\newcommand{\avrg}[1]{\left\langle #1 \right\rangle}
\newcommand{\nelta}{\bar{\delta}}
\newcommand{\bra}[1]{\left\langle #1\right|}
\newcommand{\ket}[1]{\left| #1 \right\rangle}
\newcommand{\sbra}[1]{\langle #1|}
\newcommand{\sket}[1]{| #1 \rangle}
\newcommand{\bek}[3]{\left\langle #1 \right| #2 \left| #3 \right\rangle}
\newcommand{\sbek}[3]{\langle #1 | #2 | #3 \rangle}
\newcommand{\braket}[2]{\left\langle #1 \middle| #2 \right\rangle}
\newcommand{\ketbra}[2]{\left| #1 \middle\rangle \middle\langle #2  \right|}
\newcommand{\sbraket}[2]{\langle #1 | #2 \rangle}
\newcommand{\sketbra}[2]{| #1 \rangle  \langle #2 |}
\newcommand{\norm}[1]{\left\lVert#1\right\rVert}
\newcommand{\snorm}[1]{\lVert#1\rVert}
\newcommand{\bvec}[1]{\boldsymbol{\mathsf{#1}}}
\newcommand{\bcov}[1]{\boldsymbol{#1}}
\newcommand{\bdua}[1]{\boldsymbol{\check{#1}}}
%\newcommand{\bdov}[1]{\boldsymbol{\breve{#1}}}
\newcommand{\bdov}[1]{\breve{#1}}
%\newcommand{\bten}[1]{\boldsymbol{\mathfrak{#1}}}
\newcommand{\bten}[1]{\boldsymbol{\mathfrak{#1}}}
\newcommand{\forany}{\tilde{\forall}}
\newcommand{\qed}{$\overset{\circ}{.}\;$}
\newcommand{\bigo}{\mathcal{O}}
\newcommand{\spn}{\mathrm{spn}\,}
\newcommand{\krn}{\mathrm{krn}\,}
\newcommand{\biy}{\leftrightarrow}
\newcommand{\ugh}{\mathbin{\textlinb{\BPwheel}}}

\newcommand\bigeye{\ensurestackMath{\stackinset{c}{}{c}{-.3pt}%
  {\bullet}{\scriptstyle\bigcirc}}}
\newcommand\eye{\scalerel*{\bigeye}{x}}
%\newcommand*{\fisheye}{%
%    \mathbin{%
%        \ooalign{$\circledcirc$\cr\hidewidth$\bullet$\hidewidth}%
%    }%
%}

% double angle-bracket
% https://tex.stackexchange.com/questions/79657/how-to-get-double-angle-bracket-without-using-mnsymbol-package
\makeatletter
\newsavebox{\@brx}
\newcommand{\llangle}[1][]{\savebox{\@brx}{\(\m@th{#1\langle}\)}%
  \mathopen{\copy\@brx\mkern2mu\kern-0.9\wd\@brx\usebox{\@brx}}}
\newcommand{\rrangle}[1][]{\savebox{\@brx}{\(\m@th{#1\rangle}\)}%
  \mathclose{\copy\@brx\mkern2mu\kern-0.9\wd\@brx\usebox{\@brx}}}
\makeatother

% https://tex.stackexchange.com/questions/297362/bar-through-letter-in-mathmode
%\usepackage{calc}
\newsavebox\CBox
\newcommand\hcancel[2][0.5pt]{%
  \ifmmode\sbox\CBox{$#2$}\else\sbox\CBox{#2}\fi%
  \makebox[0pt][l]{\usebox\CBox}%  
  \rule[0.77\ht\CBox-#1/2]{\wd\CBox}{#1}}
\newcommand{\cb}{\hcancel{b}}
\newcommand{\cd}{\hcancel{d}}
\newcommand{\ct}{\hcancel{\partial}}
%\newcommand{\supp}{\mathsf{supp}}
\newcommand{\supp}{\mathrm{supp}}
\newcommand{\csupp}{\overline{\mathrm{supp}}}
\newcommand{\supr}[2]{\underset{#1}{\text{supr}}}

\begin{document}

% Use the \preprint command to place your local institutional report
% number in the upper righthand corner of the title page in preprint mode.
% Multiple \preprint commands are allowed.
% Use the 'preprintnumbers' class option to override journal defaults
% to display numbers if necessary
%\preprint{}

%Title of paper
\title{
M\'etodos Matem\'aticos de la F\'isica II\\
Partial Differential Equations
}

% repeat the \author .. \affiliation  etc. as needed
% \email, \thanks, \homepage, \altaffiliation all apply to the current
% author. Explanatory text should go in the []'s, actual e-mail
% address or url should go in the {}'s for \email and \homepage.
% Please use the appropriate macro foreach each type of information

% \affiliation command applies to all authors since the last
% \affiliation command. The \affiliation command should follow the
% other information
% \affiliation can be followed by \email, \homepage, \thanks as well.
\author{Juan I. Perotti}
\email[]{juan.perotti@unc.edu.ar}
%\homepage[]{Your web page}
%\thanks{}
%\altaffiliation{}
%\affiliation{}
\affiliation{Instituto de F\'isica Enrique Gaviola (IFEG-CONICET), Ciudad Universitaria, 5000 C\'ordoba, Argentina}
\affiliation{Facultad de Matem\'atica, Astronom\'ia, F\'isica y Computaci\'on, Universidad Nacional de C\'ordoba, Ciudad Universitaria, 5000 C\'ordoba, Argentina}

% \author{Colaborador}
% \affiliation{Instituto de F\'isica Enrique Gaviola (IFEG-CONICET), Ciudad Universitaria, 5000 C\'ordoba, Argentina}
% \affiliation{Facultad de Matem\'atica, Astronom\'ia, F\'isica y Computaci\'on, Universidad Nacional de C\'ordoba, Ciudad Universitaria, 5000 C\'ordoba, Argentina}
% \email[E-mail: ]{xxx.yyy@unc.edu.ar }

%Collaboration name if desired (requires use of superscriptaddress
%option in \documentclass). \noaffiliation is required (may also be
%used with the \author command).
%\collaboration can be followed by \email, \homepage, \thanks as well.
%\collaboration{Claudio J. Tessone}
%\noaffiliation

\date{\today}

%\begin{abstract}
%Bla bla bla... resumen...
%\end{abstract}

% insert suggested keywords - APS authors don't need to do this
%\keywords{}

%\maketitle must follow title, authors, abstract, and keywords
\maketitle
\onecolumngrid

% https://chatgpt.com/share/6a21decf-f3e8-83e9-bd6c-97000a10761c

\section{Classification}

Consider an $n$-dimensional manifold $D$.
Consider functions
$$
a^{ij}:D\times \RR\times \RR^n \ni (x,v,p) \to a^{ij}(x,v,p) \in \RR
$$
for $i,j=1,2,...,n$ and
$$
f:D\times \RR\times \RR^n \ni (x,v,p) \to f(x,v,p) \in \RR.
$$
Then, we can defime a so called semi-linear Partial Differential Equations (PDEs) as
\begin{equation}
\label{eq:X1}
a^{ij}(x,u(x),(\nabla u)(x))
(
\nabla_i
\nabla_j
u)
(x)
+
f(x,u(x),(\nabla u)(x))
=
0
\end{equation}
where 
$$
u:D\ni x \to u(x)\in \RR
$$ 
is a solution and $(\nabla u)(x)\in \RR^n$ is the tangent vector of $D$ at $x$.

For each $x\in D$ and solution $u$, the matrix $a$ of coefficients
$$
a^{ij}(x,u(x),(\nabla u)(x))
\in \RR
$$
is assumed to be symmetric.
For a given $x\in D$ and solution $u$, the PDE is:
\begin{itemize}
\item {\em elliptic} at $(x,u(x),(\nabla u)(x))$ if the eigenvalues of the matrix $a$ are non-zero and of the same sign,

\item {\em hyperbolic} at $(x,u(x),(\nabla u)(x))$ if the eigenvalues of the matrix $a$ are non-zero but of different signs, and

\item {\em parabolic} at $(x,u(x),(\nabla u)(x))$ if at least one of the eigenvalues of the matrix $a$ is zero and the non-zero ones have the same sign.
\end{itemize}

\section{Boundary Conditions}

To obtain a unique solution of a PDE, one typically supplements the equation with appropriate boundary and/or initial conditions.
The most common boundary conditions for second-order PDEs are summarized by
\begin{equation}
\label{eq:X2}
a(x)u(x) + b(x)(\hat{n}\cdot (\nabla u))(x)
=
g(x)
\end{equation}
where $\hat{n}(x)$ is the unit outward normal vector to boundary $\partial D$ of $D$ at $x\in \partial D$.
The boundary condition is of type:
\begin{itemize}
\item {\em Dirichlet} if $a\neq 0$ and $b=0$,
\item {\em Neumann} if $a=0$ and $b\neq 0$,
\item {\em mixed} if $\partial D=\Gamma_D\cup\Gamma_N$ with
      $\Gamma_D\cap\Gamma_N=\varnothing$,
      and the condition is Dirichlet on $\Gamma_D$
      and Neumann on $\Gamma_N$,
\item {\em Robin} if $a$ and $b$ are non-zero.
\end{itemize}
Also, the boundary condition is homogeneous if $g=0$, and non-homogeneous otherwise.

For some PDEs, pure Neumann boundary conditions do not determine the solution uniquely. 
If $u$ is a solution, then $u+C$ is also a solution for any constant $C$.
This is because the Neumann condition 
$$
\frac{\partial u}{\partial \hat{n}}(x)
=
\frac{g(x)}{b(x)}
=:
h(x)
$$
is defined modulo an aditive constant and because, due to Gauss theorem,
$$
\int_{\partial D} h(x)\,dS(x)
=
\int_{\partial D} (\hat{n}\cdot (\nabla u))(x)\,dS(x)
=
\int_{D} (\Delta u)(x)\,dS(x)
=
\int_{D} v(x)\,dV(x)
$$
for some function $v:D\to \RR$ that is determined by a solution $u$.
In other words, the function $h$ should be compatible with one of the  possible solutions.

\section{(Unusual) weak formulation}

% https://chatgpt.com/share/6a22a9be-6d98-83e9-a081-7bfe6b025c72

Because of Eq.~\ref{eq:X1}, equation
\begin{equation}
\label{eq:X3}
\int_D
\bigg(
a^{ij}(x,u(x),(\nabla u)(x))
\,
(\nabla_i \nabla_j u)(x)
+
f(x,u(x),(\nabla u)(x))
\bigg)\,
\varphi(x)\,
dV(x)
=
0
\end{equation}
holds for every test function $\varphi\in C_c^{\infty}(D)$ if $u$ is a solution.
Conversely, any $u$ satisfying Eq.~\ref{eq:X3} defines a distributional (or weak) solution of the PDE in the interior of $D$. 
Since the test functions belong to $C_c^{\infty}(D)$, the boundary is not involved and no boundary condition is encoded. 
Any Dirichlet, Neumann, or Robin condition must be imposed separately.
The existence of distributional solutions, enlarge the space of possible solutions.
For instance, solutions can have distributional like boundary conditions.

\section{Divergence form}

% https://chatgpt.com/share/6a22a9be-6d98-83e9-a081-7bfe6b025c72

The weak formulation can be rewritten in a more convenient way, if Eq.~\ref{eq:X1} is rewriten in the so called divergence form.
To do this, consider an arbitrary solution $u$ of Eq.~\ref{eq:X1}, and let us define the functions $A^{ij}:D\ni x\to A^{ij}(x)\RR$ by
$$
A^{ij}(x)
:=
a^{ij}(x,u(x),(\nabla u)(x)),
$$
so we can write
$$
a^{ij}(x,u(x),(\nabla u)(x))\,
(\nabla_i \nabla_j u)
(x)
%=
%A^{ij}(x)
%(
%\nabla_i
%\nabla_j
%u)
%(x)
=
(A^{ij}
\nabla_i
\nabla_j
u)
(x)
.
$$
Next, consider the identity obtained by Leibniz rule
$$
\nabla_i
\big(
A^{ij}
\nabla_j
u
\big)
=
\big(
\nabla_i
A^{ij}
\big)
\nabla_j
u
+
A^{ij}
\nabla_i
\nabla_j
u
$$
which can be rearranged into
$$
A^{ij}
\nabla_i
\nabla_j
u
=
\nabla_i
\big(
A^{ij}
\nabla_j
u
\big)
-
\big(
\nabla_i
A^{ij}
\big)
\nabla_j
u
.
$$
In this way, after defining the function $F:D\ni x\to F(x)=f(x,u(x),(\nabla u)(x))\in \RR$, Eq.~\ref{eq:X1} can be rewritten as
\begin{eqnarray}
\label{eq:X7}
0
&=&
A^{ij}
\nabla_i
\nabla_j
u
+
F
\notag
\\
&=&
\nabla_i
\big(
A^{ij}
\nabla_j
u
\big)
+
F
-
\big(
\nabla_i
A^{ij}
\big)
\nabla_j
u
\notag
\\
-
\nabla_i
\big(
A^{ij}
\nabla_j
u
\big)
+
\tilde{F}
&=&
0
%\notag
\end{eqnarray}
where
\begin{equation}
\label{eq:X8}
\tilde{F}
=
\big(
\nabla_i
A^{ij}
\big)
\nabla_j
u
-
F
.
\end{equation}
Equation~\ref{eq:X7} is called the divergence form of Eq.~\ref{eq:X1}.

Sometimes the so called divergence free condition
$$
\nabla_i A^{ij} = 0
$$
holds, implying ...
{\color{red}[TODO]}

\section{Weak formulation revisited}

% https://chatgpt.com/share/6a22c3f4-6860-83e9-a018-452f279a772a

Now we revisit the weak formulation but using the divergence form of the PDE.
Namely, we multiply Eq.~\ref{eq:X7} by a test function $\varphi$ and integrate it to obtain
$$
-
\int_D
\nabla_i
\big(
A^{ij}
\nabla_j
u
\big)\,
\varphi\,
dV
+
\int_D
\tilde{F}\,
\varphi\,
dV
=
0
.
$$
Hence, by applying the divergence theorem, this equation can be rewriten as
$$
\int_D
A^{ij}\,
(\nabla_i \varphi)
(\nabla_j u)
\,
dV
-
\int_{\partial D}
\varphi\,
\big(
A^{ij}
\nabla_j
u
\big)\,
\hat{n}_i\,
dS
+
\int_D
\tilde{F}\,
\varphi\,
dV
=
0
.
$$
{\color{red}[INCORRECT... se ChatGPT]}
Notice what happened:
\begin{itemize}
\item the second derivatives of $u$ disappeared,
\item only first derivatives of $u$ remain, and
\item boundary terms appeared.
\end{itemize}
This is convenient because the weak formulation can then be defined for functions $u$ that have only weak first derivatives rather than the original Eq.~\ref{eq:X1} that requires $u$ to have well defined second derivatives in the strong sense.

%\appendix
%
%\section{
%Compactness in infinite dimension
%\label{app:A}
%}
%
%% https://chatgpt.com/share/6a0c5541-7498-83e9-a94f-8ca91aa5e791
%
%The compactness of a subset $K$ of a metric space $X$ only makes sense if $X$ is finite-dimensional?
%No. Compactness is a fundamental topological notion and makes perfect sense in {\em any} metric space, including infinite-dimensional ones.
%For a subset $K \subset X$ of a metric space $(X,d)$, compactness means: 
%\begin{quote}
%Every open cover of $K$ admits a finite subcover.
%\end{quote}
%In metric spaces, this is equivalent to: 
%\begin{quote}
%Every sequence in $K$ has a convergent subsequence whose limit lies in $K$.
%\end{quote}
%This definition does not depend on dimension at all.
%What {\em does} change drastically between finite- and infinite-dimensional spaces is {\bf which sets are compact}.
%
%\subsection{Finite-dimensional case}
%
%In $\mathbb{R}^n$, the Heine–Borel theorem says:
%$$
%K \subset \mathbb{R}^n
%\text{ is compact }
%\iff
%K \text{ is closed and bounded.}
%$$
%This is a very special property of finite-dimensional normed spaces.
%
%\subsection{Infinite-dimensional case}
%
%In infinite-dimensional normed spaces, closed and bounded sets are {\em usually not compact}.
%For example, consider the Hilbert space
%$$
%\ell^2
%=
%\left\{
%x=(x_1,x_2,\dots):
%\sum_{n=1}^\infty |x_n|^2 < \infty
%\right\}
%.
%$$
%Its closed unit ball
%$$
%B=\{x\in \ell^2:|x|\le 1\}
%$$
%is closed and bounded, but {\bf not compact}.
%A standard proof uses the canonical basis vectors
%$$
%e_1=(1,0,0,\dots),\quad
%e_2=(0,1,0,\dots),\quad \dots
%$$
%which satisfy
%$$
%|e_n-e_m|=\sqrt{2}
%\quad (n\neq m).
%$$
%Hence no subsequence can be Cauchy, so no subsequence converges. Therefore $B$ is not compact.
%
%\subsection{What replaces Heine–Borel?}
%
%In metric spaces, compactness is closely related to {\bf total boundedness}.
%
%For metric spaces:
%
%$$
%K \text{ compact}
%\iff
%K \text{ complete and totally bounded}.
%$$
%
%Total boundedness is stronger than ordinary boundedness.
%A set is totally bounded if:
%\begin{quote}
%For every $\varepsilon>0$, it can be covered by finitely many balls of radius $\varepsilon$.
%\end{quote}
%In finite dimensions, boundedness already implies total boundedness. In infinite dimensions, it generally does not.
%
%\subsection{Intuition}
%
%Infinite-dimensional balls contain ``too many independent directions.''
%The sequence $e_n$ in $\ell^2$ keeps jumping into new orthogonal directions and never accumulates anywhere.
%That phenomenon cannot happen inside bounded subsets of finite-dimensional spaces.
%So compactness absolutely makes sense in infinite-dimensional spaces — it simply becomes much more restrictive.


\end{document}
%
% ****** End of file apstemplate.tex ******


